\documentclass{article}
\usepackage{amsmath,amssymb}
\newcommand{\uh}{u_h}
\newcommand{\qh}{\boldsymbol q_h}
\newcommand{\bh}{\boldsymbol\beta_h}
\newcommand{\wideh}{\widehat u_h}
\begin{document}
\title{Stationary Advection--Diffusion--Reaction HDG}
\maketitle

We solve
\[
 \nabla\!\cdot(\bh\uh+\qh)+r_h\uh=f_h,
 \qquad \qh=-\kappa\nabla\uh,
 \qquad \wideh=g_D\quad\hbox{on }\partial\Omega .
\]
Normals are outward from each element.  The total numerical normal flux is
\[
 \widehat{\boldsymbol F}_h\!\cdot\boldsymbol n
 =\qh\!\cdot\boldsymbol n+(\bh\!\cdot\boldsymbol n)\wideh
  +(\tau_{\rm adv}+\tau_{\rm diff})(\uh-\wideh).
\]
Set $\tau=\tau_{\rm adv}+\tau_{\rm diff}$ and
$\gamma=\tau-\bh\!\cdot\boldsymbol n$.  With scalar test $w$, flux test
$\boldsymbol v$, and trace test $\mu$, the element equations are
\begin{align*}
 (\kappa^{-1}\qh,\boldsymbol v)_K-(\uh,\nabla\!\cdot\boldsymbol v)_K
       +\langle\wideh,\boldsymbol v\!\cdot\boldsymbol n\rangle_{\partial K}&=0,\\
 -(\bh\uh+\qh,\nabla w)_K+(r_h\uh,w)_K
       +\langle\widehat{\boldsymbol F}_h\!\cdot\boldsymbol n,w\rangle_{\partial K}
       &=(f_h,w)_K.
\end{align*}
Thus the scalar--scalar local block contains reaction, negative volume
advection, and the boundary mass weighted by $\tau$.  The trace-to-scalar
local right-hand block is weighted by $\gamma$; the trace-to-flux blocks are
the outward-normal face moments.  If $x_K=(u_K,q_{x,K},q_{y,K})$ and $\lambda_K$
contains the three oriented element traces, this has the form
\[
 A_Kx_K=f_K+B_K\lambda_K,
 \qquad x_K=A_K^{-1}(f_K+B_K\lambda_K).
\]

The transmission condition on every interior edge $F$ is
\[
 \sum_{K:\,F\subset\partial K}
 \left\langle
   \qh\!\cdot\boldsymbol n+\tau\uh-\gamma\wideh,\mu
 \right\rangle_F=0.
\]
Consequently each element incidence contributes its own Schur block and its
own negative $\gamma_{K,F}$ trace-mass block.  These blocks must be emitted in
the element/local-face loop.  In particular, two sides of one global edge are
not replaced by an averaged stabilization and are not represented by a
hard-coded factor of two.  This is required for element- and face-dependent
stabilization.

The default advective stabilization is upwind,
$\tau_{\rm adv}=|\bh\!\cdot\boldsymbol n|$.  For positive constant scalar
diffusion the production default is
\[
 \tau_{{\rm diff},K,F}=\gamma_d\frac{\kappa}{L_\Omega},
 \qquad L_\Omega=\frac{2|\Omega|}{|\partial\Omega|},\qquad \gamma_d=1.
\]
A positive user-supplied physical $L_\Omega$ takes precedence over the
automatic geometry value.  Explicit positive values may be supplied per
element side.  Source, reaction,
and velocity are sampled at the quadrature needed by the solve space; their DG
spaces need only share the mesh and need not equal the unknown space.

The legacy $C_\tau(p+1)^2\kappa/h_{K,F}$ rule remains an explicit comparison
mode.  It is not assumed in the postprocessing estimate.  Both recoveries
below reproduce the expected $O(h^{p+1})$ total-flux and $O(h^{p+2})$ primal
rates in the smooth manufactured test with the positive $h$-independent
default.  With the inverse-$h$ scaling, the same test instead exhibits
approximately $O(h^p)$ total-flux and $O(h^{p+1})$ primal behavior.

After trace solution and local reconstruction, we first reconstruct a
conservative total flux $\boldsymbol q_h^{T,*}$.  The default method seeks it
in $[P_{p+1}(K)]^2$, matches normal moments of
$\widehat{\boldsymbol F}_h\!\cdot\boldsymbol n$ against $P_{p+1}(F)$ and
component moments of $\qh+\bh\uh$ against $[P_{p-1}(K)]^2$, and minimizes the
remaining elementwise $L^2$ distance to $\qh+\bh\uh$.

An experimental alternative seeks the same total flux in the Raviart--Thomas
space
\[
 RT_p(K)=[P_p(K)]^2+\boldsymbol x P_p(K).
\]
The implementation uses $[P_p(K)]^2$ plus the nonredundant
$\boldsymbol x\widetilde P_p(K)$ homogeneous-degree-$p$ complement.  Its
$3(p+1)$ normal moments against $P_p(F)$ and $p(p+1)$ component moments
against $[P_{p-1}(K)]^2$ give exactly
$(p+1)(p+3)=\dim RT_p(K)$ equations, so this reconstruction is unisolvent and
has no $L^2$ minimization.  The interior targets are evaluated directly from
$\qh+\bh\uh$ on the postprocessing quadrature, rather than from a preliminary
degree-$p$ projection.  Both alternatives are stored componentwise in the
existing degree-$p+1$ DG space.

The primal reconstruction is not the diffusion-only recovery.
Instead, with $k^*=p+1$, it finds
$(u_h^*,\boldsymbol q_h^*,\phi_h^*)$ in
$P_{k^*}(K)\times[P_{k^*}(K)]^2\times P_{k^*}(\partial K)$ such that
\begin{align*}
 (\kappa^{-1}\boldsymbol q_h^*,\boldsymbol v)_K
 -(u_h^*,\nabla\!\cdot\boldsymbol v)_K
 +\langle\phi_h^*,\boldsymbol v\!\cdot\boldsymbol n\rangle_{\partial K}&=0,\\
 -(\boldsymbol q_h^*+\bh u_h^*,\nabla w)_K
 +\langle\widehat F_h^*,w\rangle_{\partial K}
 &=(\nabla\!\cdot\boldsymbol q_h^{T,*},w)_K,\\
 \langle\widehat F_h^*,\mu\rangle_{\partial K}
 &=\langle\boldsymbol q_h^{T,*}\!\cdot\boldsymbol n,\mu\rangle_{\partial K},\\
 (u_h^*,1)_K&=(\uh,1)_K,
\end{align*}
where
\[
 \widehat F_h^*=\boldsymbol q_h^*\!\cdot\boldsymbol n
 +(\bh\!\cdot\boldsymbol n)\phi_h^*
 +\tau(u_h^*-\phi_h^*).
\]
The implementation evaluates the divergence right-hand side by integration by
parts from $\boldsymbol q_h^{T,*}$ and adds one scalar Neumann multiplier for
the mean constraint.  Thus ADR primal postprocessing always performs the total
flux reconstruction first.  The available output modes remain
\texttt{none}, \texttt{primal}, \texttt{flux}, and \texttt{both}; in
\texttt{primal} mode the prerequisite total flux is internal.

The NumPy implementation is the dense local reference.  The multithreaded
Numba implementation fuses element construction, local elimination, and COO
emission with \texttt{prange}.  The raw-CUDA implementation repeats the same
incidence algebra on the device, converts device COO to CSR, solves through
AMGX, and reconstructs element unknowns on the device.
\end{document}
